%% notes-tex/010-index-defect/sections/4-negatives.tex

\section{What the model can still be used for\secstatus{todo}}
\label{sec:negatives}

Every confident prediction in Section~\ref{sec:panel} is negative. That suggests
using the model to nominate negative trigenic interactions rather than positive
ones, and it is worth testing rather than assuming, because the label is already
skewed negative: 59.4 percent of the 376{,}732 records have $\tau$ below zero and
the mean is $-0.0080$. A squared-error fit shrinks toward the mean, so a model
that had learned nothing beyond it would also emit mostly negative predictions.
That most predictions are negative is therefore not evidence of anything.

\subsection{The measurement that is evidence\secstatus{todo}}

Retrieval. Among the $K$ records a model calls most negative, what fraction truly
are strong negatives, and does that beat both the base rate and the additive
null.

%% SOURCE: experiments/010-kuzmin-tmi/scripts/negative_interaction_retrieval.py
%% -> results/negative_interaction_retrieval.csv, held-out test, 37,673 records
\begin{table}[htbp]
\centering
\begin{tabular}{lrrr}
\hline
Model & $\tau < -0.08$ & $\tau < -0.10$ & $\tau < -0.20$ \\
 & base 8.3\% & base 5.1\% & base 0.88\% \\
\hline
B1 additive ridge & 0.560 & 0.520 & 0.240 \\
B5 nonlinear MLP & 0.700 & 0.670 & 0.400 \\
CGT M01 & 0.670 & 0.630 & 0.400 \\
CGT M02 & 0.750 & 0.700 & 0.420 \\
CGT M03 & 0.710 & 0.670 & 0.450 \\
\hline
\end{tabular}
\caption{Precision at $K = 100$: the fraction of the 100 most negative
predictions whose true $\tau$ clears each threshold. At the strictest threshold
the transformer converts a 0.88 percent base rate into 42 to 45 percent, a 48 to
51 fold enrichment.}
\label{tab:retrieval}
\end{table}

The transformer's advantage over the additive null is proportionally larger here
than on correlation. Average precision at $\tau < -0.20$ is $0.096$ for the ridge
against $0.182$ to $0.191$ for the three checkpoints, whereas on Pearson the same
comparison was $0.400$ against $0.438$ to $0.455$. Retrieving the extreme
negative tail is what the extra capacity buys, and it is a better description of
the model's competence than its correlation is.

\subsection{It survives the disjoint split, several times weaker\secstatus{todo}}

The same measurement was run on the query-pair-disjoint split, where the leak
described in the companion document is removed. Only the baselines have been
refit there, so the transformer has no row.

At $\tau < -0.10$ the disjoint test base rate is 9.6 percent, and precision at
100 is $0.26$ for the additive ridge and $0.33$ for the nonlinear baseline,
enrichments of 2.7 and 3.4 fold against 10.1 and 13.1 on the random split.
Absolute precisions are not comparable across the two splits because their base
rates differ; enrichment is the quantity that is.

\tcfigfit[0.98]{figures/negative_interaction_retrieval.pdf}{Precision at $K$ for
retrieving records with true $\tau < -0.10$, on both splits, with the base rate
marked.}{fig:retrieval}

\subsection{What this supports\secstatus{todo}}

Using the model to nominate candidate negative trigenic interactions is
defensible on this evidence and fits what it can do better than nominating
positives. On the build list 17 of 20 triples have all three checkpoints agreeing
and every one is negative, while the single positive is not agreed.

The enrichment a transformer would retain once query doubles are disjoint is not
established, because it has not been refit there. The baselines' fall from
roughly tenfold to threefold is the honest prior for what to expect.
